


Rule (V14 section 5): a protocol claim reaches the main text only if it holds in FORMAL or TRACE\_DRIVEN evidence, or in at least two of the three link models. Seven preregistered tests share one Holm family (adjusted p about 0.0035, set by the resolution of the bootstrap); H-T2, H-T2sec and H-T3c are decided by their stated criteria. Nearest prior work is the closest item of \texttt{docs/06\_references/NOVELTY\_MATRIX\_2021PLUS.md}.

{\scriptsize\raggedright
\begin{longtable}{p{0.136\linewidth}p{0.136\linewidth}p{0.136\linewidth}p{0.136\linewidth}p{0.136\linewidth}p{0.136\linewidth}p{0.136\linewidth}}
\toprule
Hypothesis & Result & Label & Interval & Nearest prior work & Allowed wording & Placement \\
\midrule
\endhead
H-T1a (L, M, H) & SUPPORTED, SUPPORTED, SUPPORTED: D1 is within 2 points of the best baseline in each group (it is above it: best minus D1 = -0.1, -1.4, -6.7 points) & TRACE\_DRIVEN & upper limits 0.0, -0.8, -6.5 points & Nessa et al. (CNS 2020, conference); bidirectional LoRa (IEEE Commun. Lett. 2021); single-channel LoRa gateway (Sci. Rep. 2024) & On measured London loss sequences, a table-driven retry policy was not worse than six other retry rules on single-link delivery within 10 s; the replay has no crowd and its delivery follows the number of attempts that fit before the deadline. & Main text only together with the zero-spacing control \\
H-T1b & SUPPORTED: D1 minus quick retries = 6.7 points on group H & TRACE\_DRIVEN & 6.5 to 6.8 points & as above & On links with 40\% loss or more, D1 delivered 6.7 points more than quick retries on the replay; a no-hold control that retries at once delivered 54.0 against D1's 49.3, so the gain is retry timing, not a mechanism shown by the traces. & Main text with the control (the rule admits it as TRACE\_DRIVEN evidence); D1 stays a supplement item in the ns-3 reading \\
H-T1c & SUPPORTED: the provisional alert arrives 1.87 s earlier (median) in every group; extra airtime 37\% to 31\%; delivery not lower & TRACE\_DRIVEN & airtime upper limits below 40\% in all groups; delivery lower limits 0.0 to 9.0 points & no journal article from 2021 onwards found in the sources checked on provisional-then-confirmed crash alerts & On the same measured traces, a 14-byte provisional flash arrived about 1.9 s before the full alert with about a third more airtime. & Main text \\
H-T2 (timed) & SUPPORTED (FORMAL): D1's worst-case delivery time is not longer than quick retries' for k = 0 to 6; no duty breach & FORMAL (exhaustive enumeration of the frozen code; not UPPAAL) & exact values in ROUND3\_RESULTS.md & none found for UPPAAL deadline bounds of LoRa alert retries in 2021+ journals & In a single-link timing model, D1 delivers within 10 s after up to six lost uplinks where quick retries do not, and never exceeds the 1\% hourly allowance. & Main text with the tool named \\
H-T2sec & PARTLY SUPPORTED: L1 to L4b and no double acceptance across a car reboot proved on M1; attack found on M0; L6 refuted as expected; id uniqueness after a car reboot not proved (timeout) & FORMAL & not applicable & Wesemeyer et al. (EuroS\&P 2020); Dong et al. (ISSRE 2024); LoRaWAN authentication analyses (IEEE TII 2021, IJSNET 2025) & A Tamarin model of a crash-alert protocol with per-vehicle keys and an ACK bound to the vehicle: authenticity, per-vehicle ACK binding, window and seen-set properties proved; a replay after a base reboot is possible; one reboot lemma not completed. & Main text for L1 to L4b and the M0 attack (R4); the unfinished lemma in the limits \\
H-T3a & SUPPORTED: median 1750 ms earlier & MEASURED\_SOFTWARE & 1750 to 1750 ms & direction-aware telematics detection (IJLTEMAS 2026); eCall injury prediction papers & On the de-duplicated VZCrash test split, the early-exit rule decides a median of 1.75 s before the full window (a stored seed-42 decision set that R3 already used). & Main text \\
H-T3b & SUPPORTED: retraction share 10.8\% & MEASURED\_SOFTWARE & 9.7\% to 12.0\% & as above & At the 1\% false-alarm target, about 1 in 9 provisional alarms was withdrawn by the full window. & Main text \\
H-T3c & SUPPORTED: holds in 2 of 3 link models (custom, trace-driven; ns-3 NOT\_RUN) & SIMULATED and TRACE\_DRIVEN & not applicable & as above & The earlier arrival of the provisional alert holds in the custom simulator and on the measured traces; ns-3 was not run for this alert. & Main text with ns-3 stated as NOT\_RUN \\
T3 class agreement & NHTSA frontal pulses: 3\% agreement at 50 ms, 98\% at 200 ms; VZCrash agreement 99.8\% (uninformative: nearly all class 0) & MEASURED\_SOFTWARE & VZCrash 99.7\% to 100.0\% & Stitzel et al. (OTDA, ESV); eCall injury papers & A severity class from the window truncated at 50 ms agrees with the full-window class for about 3\% of frontal NHTSA pulses; at 200 ms for 98\%. & Supplement (limit of the provisional alert) \\
T4 energy budget & no hypothesis; at +17 dBm the all-lost schedule of D1 needs about 5 J, most of it for the ESP32 and GPS being awake & ANALYTICAL & not applicable & none found (topic 13) & Under stated assumptions, a 1 F capacitor between 5 V and 3 V can supply the whole 10 s schedule of every policy. & Supplement \\
\bottomrule
\end{longtable}}

\subsection*{What the round does not support}

\begin{itemize}
\item A mechanism of D1 beyond retry timing: the zero-spacing control beats it on the traces, and the ns-3 gain did not reproduce.
\item Any statement about vehicles, ACK exchanges or crowds from the traces.
\item An implementation guarantee from the Tamarin models, and the id-uniqueness lemma after a car reboot (not proved).
\item Any severity class in the provisional alert: at 50 ms the class is not yet determined for frontal crashes.
\item Any claim that no prior work exists: use the wording of the novelty matrix (no journal article from 2021 onwards found in the sources checked; Scopus, Google Scholar and TRID not run; IEEE Xplore not run in this session).
\end{itemize}
